Label each edge by the coordinate it changes. Each of the three coordinates must change an odd number of times. At each of the six other vertices a used-edge degree is zero or two; the start has degree at most three and the destination degree one. Thus there are at most eight used edges, and the length is odd, hence at most seven.
Length three uses each coordinate once, giving $3!=6$ paths. At length five the multiplicities are $(3,1,1)$. The repeated coordinate $X$ cannot occur consecutively, so the sequence must be $XYXZX$. There are $3\cdot2=6$ such paths.
There are only four cube edges of each coordinate type, so an odd multiplicity cannot exceed three. At length seven the multiplicities are therefore $(3,3,1)$. Choose the coordinate used once in three ways. Its edge switches between two opposite square faces. An edge-simple walk within a square follows one direction and has length at most four. The six remaining steps split as $(2,4),(3,3)$ or $(4,2)$. The first split reaches the destination at the switch and is forbidden. For $(3,3)$ there are two first-face directions and the second-face direction is forced. For $(4,2)$ there are two directions for the first circuit and two for the last two-step route. Thus there are $3(2+4)=18$ seven-step paths. The total is $\boxed{6+6+18=30}$.
\begin{asy}
texpreamble("\usepackage{amsmath,amssymb}");
size(240);
defaultpen(fontsize(9pt));
pair A=(0,0),B=(3,0),C=(4,1),D=(1,1),E=(0,3),F=(3,3),G=(4,4),H=(1,4);draw(A--B--C--D--cycle);draw(E--F--G--H--cycle);draw(A--E);draw(B--F);draw(C--G);draw(D--H,dashed);dot(A);dot(G);label("$A$",A,(-1,-1));label("$C_1$",G,(1,1));
\end{asy}
